\section{The model, normalization and topologies}\label{intro:model-section}

\subsection{The finite FK law}

\begin{definition}[Finite ensemble]\label{intro:finite-law}
Let $\mathcal F_n$ consist of pairs $(M,A)$, where $M$ is a connected
planar multigraph with $n$ edges embedded in the oriented sphere and
rooted at a distinguished oriented edge, and $A\subseteq E(M)$.
Pairs are identified under orientation-preserving homeomorphisms
preserving the root and $A$. Loops and multiple edges are allowed.
Let $M^*$ be the dual map, and let $A^*$ consist of the dual edges
crossing $E(M)\setminus A$. Write $k(A)$ and $k(A^*)$ for the numbers
of connected components of the respective spanning subgraphs,
including isolated vertices. Set
\[
 \ell(M,A)=k(A)+k(A^*)-1
         =2k(A)+|A|-|V(M)|,
 \qquad
 Z_n(q)=\sum_{(M,A)\in\mathcal F_n}q^{\ell(M,A)/2}.
\]
We sample $(M_n,A_n)$ with probability
$q^{\ell(M,A)/2}/Z_n(q)$. Its graph distance $d_n$ is shortest-path
distance on $V(M_n)$ using all edges of $M_n$, each of length one.
\end{definition}

For completeness, thicken the vertices and occupied edges to a regular
neighborhood $N(A)$. It has Euler characteristic $|V(M)|-|A|$.
Every component has genus zero. If its total number of boundary
components is $b_A$, then
$|V(M)|-|A|=2k(A)-b_A$, so $b_A=2k(A)+|A|-|V(M)|$.
The complementary components correspond to the components of $A^*$;
Euler's formula on the sphere gives $b_A=k(A)+k(A^*)-1$.
This proves the displayed identities and identifies $\ell$ with the
number of FK interfaces, including configurations containing bridges
or loops.

\subsection{Flag triangles and the fixed conformal coordinate}

For each occurrence of an endpoint and a side of an original edge,
take an equilateral triangle with labels $(v,e,f)$: original vertex,
edge midpoint and face center. Glue sides by their incidence
identifications. There are $4n$ triangles. Repeated incidences are
different occurrences, even when two labels refer to the same vertex,
edge or face. In particular this construction is the incidence
subdivision of the embedded sphere and produces a topological sphere
$X_n$, also for loops and bridges.

The Euclidean metrics on the triangles define a conformal structure
away from their conical vertices. At a cone of total angle $\theta>0$,
the coordinate obtained by raising the local cone coordinate to the
power $2\pi/\theta$ extends this structure across the vertex. Give
$X_n$ normalized Euclidean face area $m_n$, so each triangle has
mass $1/(4n)$. Conditional on $X_n$, take three independent samples
$x_n^1,x_n^2,x_n^3$ from $m_n$. Since $m_n$ has no atoms they are
almost surely distinct. Genus-zero uniformization
\cite[Theorem~II.2.15]{SaintGervais2016Uniformization} and three-point
normalization give a unique orientation-preserving conformal map
\[
 \phi_n:X_n\longrightarrow\Chat,
 \qquad
 \phi_n(x_n^1)=0,\quad\phi_n(x_n^2)=1,\quad
 \phi_n(x_n^3)=\infty.
\]
This is a function of the discrete surface and the three samples.

In a triangle write $(t_v,t_e,t_f)$ for barycentric coordinates.
Include $\{t_v\ge2/3\}$ and, if the original edge belongs to $A_n$,
also include $\{t_f\le1/3\}$. Their glued union is a regular
neighborhood of the occupied primal subgraph. Its boundary components
are the prescribed polygonal interfaces. Write $\widehat\Gamma_n$
for their collection on $X_n$, taken as unrooted and unoriented loops,
and $\Gamma_n=\phi_n(\widehat\Gamma_n)$ for the canonical images.

\begin{lemma}[Flag area and a uniform root]\label{intro:root-flag}
Let $\pi_n$ send the interior of each flag triangle to its primal
vertex, with an arbitrary convention on triangle boundaries. Then
\begin{equation}\label{intro:degree-measure}
 (\pi_n)_*m_n
 =\nu_n:=\frac1{2n}\sum_{v\in V(M_n)}\deg_{M_n}(v)\delta_v.
\end{equation}
Conditional on the unrooted decorated map, the tail of a uniform
oriented root edge has law $\nu_n$, and re-rooting uniformly preserves
the finite decorated law. If $H_n:X_n\to S$ are homeomorphisms into
a compact metric sphere with
$\max_T\diam H_n(T)\to0$, then $(H_n)_*m_n$ and
$(H_n)_*\nu_n$ have the same weak subsequential limits.
\end{lemma}

\begin{proof}
A primal vertex $v$ is incident to $2\deg(v)$ flag occurrences;
a loop contributes twice to $\deg(v)$ and all four of its occurrences
are counted. Thus its mass is $2\deg(v)/(4n)$, proving
\eqref{intro:degree-measure}. There are $2n$ oriented edge
occurrences, and exactly $\deg(v)$ have tail $v$. Uniform re-rooting
preserves the weight $q^{\ell/2}$; equivalently, count rooted
occurrences before quotienting by the automorphism group. This also
proves the assertion in the presence of automorphisms.
Couple $x\sim m_n$ to $\pi_n(x)$. Their $H_n$ images lie in the
same flag image, so their distance is at most
$\max_T\diam H_n(T)$. Testing against uniformly continuous
functions on the compact sphere proves the last assertion.
\end{proof}

The area samples used for $\phi_n$ are fresh independent samples.
\Cref{intro:root-flag} describes the distribution of the forgotten
discrete root; it does not replace those three samples by exploration
marks.

\subsection{The continuum object}

The continuum surface is the ordinary $\gamma$-LQG sphere of
\citet[Definition~4.21(ii)]{DMSMating}, disintegrated at area one. Forget
its original marked points. Given the surface, sample three independent
points from its quantum-area probability measure and send them to
$0,1,\infty$ by an orientation-preserving conformal map. The area
measure is diffuse with full support, so this normalization is almost
surely defined. Write $h$ for a representative field in that
coordinate, $\mu_h$ for its area and $D_h$ for its intrinsic metric.
The metric is a continuous length metric inducing the topology of
$\Chat$. Its deterministic multiplicative normalization remains
fixed.

Given this embedded marked surface, sample the M\"obius-invariant
whole-plane nested $\mathrm{CLE}_{\kappa}$ ensemble on $\Chat$.
This conditional-law specification is part of the joint limit.
The original root and the original quantum-sphere marks are absent
from the limiting object.

\subsection{Distance graphs and loop matching}

Use the product distance
\[
 d_\times((z,w,r),(z',w',r'))
   =\rho(z,z')+\rho(w,w')+|r-r'|
\]
on $\Chat^2\times[0,\infty)$. The distance graphs
$K_n(a)$ in \eqref{intro:distance-graph} are nonempty compact sets.
Since $D_h$ is continuous on the compact sphere, its full graph is
also compact. We use the Hausdorff topology for $d_\times$.

\begin{definition}[Loop and ensemble convergence]
\label{intro:loop-topology}
Our curve distance and locally finite multiset topology follow
\citet[Section~2.2]{GwynneMillerQianCLE}, with loop orientation additionally
forgotten. They retain the order within each loop and its multiplicity.
For continuous loops $\alpha,\beta:\mathbb S^1\to\Chat$, let
\[
 d_{\mathrm{loop}}(\alpha,\beta)
 =\inf_{\sigma}\sup_{t\in\mathbb S^1}
       \rho(\alpha(t),\beta(\sigma(t))),
\]
where $\sigma$ ranges over circle homeomorphisms of either orientation.
Identify loops at zero distance. A loop collection is locally finite
in diameter if it has finitely many loops of diameter greater than
$\varepsilon$ for every $\varepsilon>0$.
An $\varepsilon$-matching of two such collections is a partial
bijection covering every loop of diameter greater than $\varepsilon$
on either side, with matched loops at distance at most $\varepsilon$;
every unmatched loop has diameter at most $\varepsilon$.
Convergence means that such matchings exist with
$\varepsilon\to0$. All collections in this paper retain nested loops.
\end{definition}

This definition also allows repetitions to be retained as separate
elements when forming a matching. The loop identification proof will
establish both local finiteness in diameter and eventual matching of
every macroscopic loop.

\subsection{Extraction conventions}

Let $d_\gamma$ denote the deterministic Hausdorff dimension of the
$\gamma$-LQG metric \citep{GwynneMillerMetric}. We use
$p=1/d_\gamma$, $\xi=\gamma p$, and
$Q=2/\gamma+\gamma/2$. The parameter $q$ always denotes the FK
weight; the inventory parameter is $b$. A crossing scale is denoted
by $B(n)$, whereas $a_n$ denotes the deterministic normalization in
\cref{intro:intrinsic-input}. The exploration root is $p_*$.

All subsequential arguments retain the same contour encoding,
continuum surface, local observations and independent area marks.
Whenever a new finite family of tests is added, refine this joint
extraction. The families used below admit countable determining
atlases: rational time cuts, rational spatial boxes, rational height
slacks and integer truncations. A diagonal refinement therefore
retains all their conclusions on a common event of probability one.
Statements involving all spatial points will be proved by explicit
uniform estimates before passing from an atlas to arbitrary points.
